% Version 0.26. The literal Bernoulli-unit and residual CorX supplier.
% Pinned source: 7cbf7a0c32fca1badf05931879a37c3f8a2b03b4.
\section{Bernoulli units and residual unramified cocycles}
\label{sec:bernoulli-residual}

\status{This section proves the prescribed Bernoulli-unit theorem
\texttt{LemmaX}, the full residual unramified cocycle theorem,
and the literal \texttt{CorX} interface.  The power-of-two conductor,
both inverse characters, the original coefficient and residue fields,
and the prime-to-$\ell$ relative-index local condition are retained.
The Bernoulli-unit argument is finite and elementary.  Its use in
cohomology rests on the Jacobi ideal relation already proved in OS1
and the finite Hilbert class field theorem stated below.}

\subsection{The exact statements and the classical boundary}

Put $G=\operatorname{Gal}(\Omega/\mathbb Q)$ for the specified algebraic
closure.  Write $I_A$ for inertia at a valuation subring $A$ of $\Omega$.
For a character $\gamma:G\to k^\times$, a cocycle means a locally
constant function $z:G\to k$ satisfying
\[
                  z(gh)=z(g)+\gamma(g)z(h).
                  \tag{BR.1}
\]
This is the source's \texttt{IsCoc}; no topology on the arbitrary field
$k$ has to be added to its statement.  Its condition
$\operatorname{UnrEv}(\ell,\gamma,z)$ says that for every $A$ there is
a subgroup $H_A\le G$ such that
\[
 \begin{gathered}
  K_A=I_A\cap\ker\gamma,\qquad
  \ell\nmid[K_A:K_A\cap H_A],\\
  z|_{K_A\cap H_A}=0.
 \end{gathered}
 \tag{BR.2}
\]
The natural-number relative index in the source is zero when the index
is infinite.  Thus the nondivisibility assertion in (BR.2) in
particular makes this index finite.  There is no requirement that
$H_A$ itself be contained in $K_A$.

Dirichlet characters are extended by zero off the units.  If $\varphi$
has level $f$ and $\xi$ has level $2^n$, let
\[
 \chi=\operatorname{prodX}(\varphi,\xi),\qquad
 \chi(a)=\varphi(a\bmod f)\xi(a\bmod 2^n).
 \tag{BR.3}
\]
The equality holds also at nonunits: a nonunit of the product level is
a nonunit at at least one factor.  We use the source normalization
\[
 B_{1,\chi}=\frac1M\sum_{a=0}^{M-1}a\chi(a)
 \quad\text{if }\chi\ne1\text{ has level }M.
 \tag{BR.4}
\]
Indeed the defining first Bernoulli polynomial is $T-\tfrac12$, and
$\sum_a\chi(a)=0$ for a nontrivial character.  Formula (BR.4) is
therefore the literal \texttt{genB 1}, including for imprimitive
characters at their specified level.
When a Dirichlet character is evaluated on $g\in G$, its argument
means the cyclotomic residue $\operatorname{cls}_M(g)$ at its stated
level $M$.  Thus
$\operatorname{redChar}(\chi)(g)=\operatorname{red}(\chi(g))$.

The data $\operatorname{XData}(\ell,d,\varphi,\xi)$ mean:
$\ell$ is prime, $d>0$, $4\mid\ell^d-1$, $f$ is odd,
$\varphi$ and $\xi$ are primitive at their respective levels,
\[
 \operatorname{ord}(\varphi)\mid\ell^d-1,\qquad
 f\,4^{v_2(\ell^d-1)}\le2^n,\qquad
 \varphi(-1)\xi(-1)=-1.
 \tag{BR.5}
\]
In particular $f>0$.  The precise Bernoulli-unit assertion is
$|B_{1,\operatorname{prodX}(\varphi,\xi)}|_\ell=1$ in the specified
$\overline{\mathbb Q}_\ell$.

\paragraph{Accepted classical input.}
For a number field $F$, the finite Hilbert class field is its maximal
abelian extension unramified at all finite places and split at its
real places; its Artin group is the ordinary ideal class group.
We use this in the form of Milne,
\href{https://www.jmilne.org/math/CourseNotes/CFT.pdf}{\emph{Class
Field Theory}}, v4.03, V, Theorems~3.5--3.6 and Example~3.9.
We also use the elementary cyclotomic identification
$\operatorname{Gal}(\mathbb Q(\mu_M)/\mathbb Q)
\simeq(\mathbb Z/M)^\times$, with $\sigma_b(\zeta)=\zeta^b$,
and extension of a finite Galois automorphism to the specified
algebraic closure; see Milne,
\href{https://www.jmilne.org/math/CourseNotes/ANT.pdf}{\emph{Algebraic
Number Theory}}, v3.08, \S6.
All transfers from these finite theorems to the actual cocycle and its
original field are given in BR4.  The needed annihilator is derived
from OS1; no general Stickelberger annihilation theorem is retained
as an additional input.

\subsection{Dyadic powers and a finite Frobenius projection}

\begin{proposition}[BR1: dyadic generation and character-sum projection]
\label{prop:br-dyadic}
Let $Q=1+2^e a$ with $e\ge2$ and $a$ odd.  For $r\ge e$, the
powers of $Q$ modulo $2^r$ are exactly the units congruent to $1$
modulo $2^e$.

Let $k$ have odd characteristic $\ell$, let $Q=\ell^d$ with
$v_2(Q-1)=e\ge2$, and let $n\ge e$.  For any finite family
$z_i,g_i\in k$ such that
$z_i=0$ or $z_i^{2^n}=1$, and $g_i^Q=g_i$, one has
\[
 \sum_i z_i g_i=0
 \quad\Longrightarrow\quad
 \sum_{\{i:z_i^{2^e}=1\}}z_i g_i=0.
 \tag{BR.6}
\]
Finally, put $n=m+e$ with $m\ge e\ge2$.  Suppose
$\eta:\mathbb Z/2^n\to k$ is multiplicative, is $1$ at $1$,
vanishes on nonunits, and
\[
                    \eta(1+2^{n-1})=-1.
                    \tag{BR.7}
\]
If $\eta(u)^{2^e}=1$, then
$u\equiv1$ or $-1\pmod{2^m}$.
\end{proposition}

\begin{proof}
Repeated squaring gives
\[
 Q^{2^j}=1+2^{e+j}a_j,\qquad a_j\text{ odd},\qquad
 a_{j+1}=a_j+2^{e+j-1}a_j^2.
 \tag{BR.8}
\]
The last summand is even, which proves the assertion about $a_j$
inductively.  Thus $Q$ has order exactly $2^{r-e}$ modulo $2^r$
when $r>e$, and order one when $r=e$.  Its powers lie in the
subgroup $1+2^e\mathbb Z/2^r$, whose size is $2^{r-e}$.
They consequently exhaust that subgroup.

Raise the assumed zero sum in (BR.6) to $Q^j$.  Frobenius is a
ring homomorphism and $g_i^{Q^j}=g_i$, so
$\sum_i z_i^{Q^j}g_i=0$.  By the first assertion, every exponent
$1+2^e t$ occurs modulo $2^n$ among the $Q^j$.  Zero entries cause
no problem, since all these exponents are positive.  Summing for
$0\le t<2^{n-e}$ gives
\[
 \begin{split}
 0&=\sum_i g_i z_i
       \sum_{t=0}^{2^{n-e}-1}(z_i^{2^e})^t\\
  &=2^{n-e}\sum_{\{i:z_i^{2^e}=1\}}z_i g_i.
 \end{split}
 \tag{BR.9}
\]
For a nonzero $z_i$ not satisfying $z_i^{2^e}=1$, the geometric
sum vanishes because its ratio has $2^{n-e}$th power one.
For $z_i=0$ the prefactor makes the term zero.  The scalar
$2^{n-e}$ is invertible in $k$, proving (BR.6).

Apply (BR.8) to $5=1+4$.
It generates all units congruent to $1$ modulo $4$ at level $2^n$.
Every odd $u$ therefore has a unique expression
$u=(-1)^\epsilon5^j$, with $j$ modulo $2^{n-2}$.
Moreover
$5^{2^{n-3}}=1+2^{n-1}$ modulo $2^n$.
By (BR.7), $\eta(5)$ has order exactly $2^{n-2}$:
its $2^{n-2}$th power is one and its $2^{n-3}$rd power is $-1$.
Since $\eta(-1)^2=1$, the condition
$\eta(u)^{2^e}=1$ implies
$2^{n-2}\mid j2^e$, or $2^{m-2}\mid j$.
The powers of $5$ with this property are $1$ modulo $2^m$,
again by (BR.8).  This proves the last assertion.
Nonunits cannot occur because their $\eta$-value is zero.
\end{proof}

\subsection{An integral partial Bernoulli sum}

\begin{proposition}[BR2: partial sums without dividing by the odd conductor]
\label{prop:br-partial}
Let $F$ be a field of characteristic zero, let $(f,N)=1$ with
$f,N>0$, and let $w:\mathbb N\to F$ be periodic modulo $f$.
Suppose a subring $B\subset F$ contains all $w(a)$ and $N^{-1}$.
There are $g:\mathbb Z/N\to B$ and $c\in B$ such that, for
$ft\le N$,
\[
 g(u)-g(u-ft)=tc-
           \begin{cases}w(\langle u\rangle_N),&
                          \langle u\rangle_N<ft,\\0,&
                          \langle u\rangle_N\ge ft,
           \end{cases}
 \tag{BR.10}
\]
where $0\le\langle u\rangle_N<N$.
For every $\xi:\mathbb Z/N\to F$ of sum zero,
\[
 \sum_{a=0}^{fN-1}a\,w(a)\xi(a)
             =fN\sum_{u\bmod N}\xi(u)g(u).
 \tag{BR.11}
\]
If $w(a)=\varphi(a\bmod f)$, $\xi$ is a nontrivial character
of level $N$, and $\varphi(-1)\xi(-1)=-1$, then the last sum is
$B_{1,\operatorname{prodX}(\varphi,\xi)}$.
\end{proposition}

\begin{proof}
Define, initially in $F$,
\[
 \begin{gathered}
 S(u)=\sum_{\substack{0\le a<fN\\a\equiv u\ (N)}}w(a),
 \qquad
 A(u)=\sum_{\substack{0\le a<fN\\a\equiv u\ (N)}}a\,w(a),\\
 c=S(0)/N,\qquad g(u)=\frac{A(u)-A(0)}{fN}.
 \end{gathered}
 \tag{BR.12}
\]
The $f$ integers in each residue class $u$ modulo $N$ run through
all residues modulo $f$, so $S(u)$ is independent of $u$.
This also follows by shifting by $f$, since $f$ generates
$\mathbb Z/N$ additively.

Here is the shift identity, including the terms crossing the endpoints.
Put $M=fN$ and
$\psi_u(a)=\mathbf1_{a\equiv u\ (N)}w(a)$.
For any function $h$ and any nonnegative $s$,
\[
 \sum_{a=0}^{M-1}h(a+s)
 =\sum_{a=0}^{M-1}h(a)
   +\sum_{a=0}^{s-1}\bigl(h(M+a)-h(a)\bigr).
 \tag{BR.13}
\]
It follows by writing the sum from $0$ to $M+s-1$ in its two
possible ways.  With $s=ft$ and $h(a)=a\psi_u(a)$, periodicity gives
$\psi_u(a+s)=\psi_{u-s}(a)$ and $\psi_u(M+a)=\psi_u(a)$.
Consequently (BR.13) becomes
\[
 A(u-s)+sS(u-s)
          =A(u)+M\sum_{a=0}^{s-1}\psi_u(a).
 \tag{BR.14}
\]
Subtract and divide by $M$ to obtain, for every $t\ge0$,
\[
 g(u)-g(u-ft)=tc-
            \sum_{\substack{0\le a<ft\\a\equiv u\ (N)}}w(a).
 \tag{BR.15}
\]
Choose $t$ with $ft=u$ modulo $N$.  Since $g(0)=0$, this formula
shows that $g(u)$ belongs to $B$.  It is only a multiple of $c$
minus a sum of the $w(a)$; division by $f$ has disappeared.
Also $c\in B$ by its definition.  When $ft\le N$, the final sum
in (BR.15) has exactly the possible term described in (BR.10).

Finally, grouping the left side of (BR.11) by residues gives
$\sum_u\xi(u)A(u)$.
The constant $A(0)$ contributes zero because $\xi$ has sum zero;
this proves (BR.11).
For the character assertion, the product character is odd and
therefore nontrivial; (BR.3)--(BR.4) identify its Bernoulli number
with the left side divided by $fN$.  This proves the asserted
identity at the original product level.
\end{proof}

\begin{proposition}[BR3: the prescribed Bernoulli unit]
\label{prop:br-lemma-x}
For every $\operatorname{XData}(\ell,d,\varphi,\xi)$ with values in
the specified $\overline{\mathbb Q}_\ell$,
\[
             |B_{1,\operatorname{prodX}(\varphi,\xi)}|_\ell=1.
             \tag{BR.16}
\]
The bound in (BR.5) is used with no enlargement or subsequent choice
of $n$.  In particular this proves the entire source
\texttt{LemmaX} interface.
\end{proposition}

\begin{proof}
Put $Q=\ell^d$ and $e=v_2(Q-1)$.  The hypotheses give $e\ge2$
and $\ell\ne2$.  Write $n=m+e$.  Since
$f\,2^{2e}\le2^n$ and $f\ge1$, one has $m\ge e$ and
\[
                   s=f2^e,\qquad 1<s\le2^m.
                   \tag{BR.17}
\]
Let $\mathcal A=\{x:|x|_\ell\le1\}$ in
$\overline{\mathbb Q}_\ell$, and let
$\mathfrak m=\{x:|x|_\ell<1\}$.  The ultrametric inequality makes
$\mathcal A$ a ring and $\mathfrak m$ an ideal.  Each element of
norm one has its inverse in $\mathcal A$, so $\mathfrak m$ is
maximal.  The field $\kappa=\mathcal A/\mathfrak m$ has
characteristic $\ell$, since $|\ell|_\ell<1$.
This auxiliary residue field is used only to test (BR.16).
It does not replace the given field $k$ in BR4--BR7.

Every value of $\varphi$ is either zero or a root of unity of
order dividing $Q-1$.  Thus it is integral and its residue is
fixed by $x\mapsto x^Q$.  Define the subring
\[
 B=\{x\in\mathcal A:\overline{x}^{\,Q}=\overline{x}\}.
 \tag{BR.18}
\]
It is a subring because the $Q$th-power map on $\kappa$ is an
iterated Frobenius.  The element $(2^n)^{-1}$ belongs to $B$:
it has norm one and its residue is an invertible prime-field element.
Apply BR2 with $N=2^n$ and $w(a)=\varphi(a\bmod f)$.
It supplies $g(u),c\in B$ and
\[
 B_{1,\operatorname{prodX}(\varphi,\xi)}
                         =\sum_{u\bmod 2^n}\xi(u)g(u).
 \tag{BR.19}
\]
The character $\xi$ is nontrivial, as follows also from its
primitivity and $n\ge4$.

For completeness, primitivity gives the exact half-period value
\[
                         \xi(1+2^{n-1})=-1.
                         \tag{BR.20}
\]
Indeed $h=1+2^{n-1}$ has order two.  The kernel of
$(\mathbb Z/2^n)^\times\to(\mathbb Z/2^{n-1})^\times$
is exactly $\{1,h\}$, as is seen from the two representatives
congruent to $1$ modulo $2^{n-1}$.
If $\xi(h)=1$, the character factors through level $2^{n-1}$,
contradicting primitivity.  In a characteristic-zero field its
value is therefore $-1$.

All values of $\xi$ are integral, so (BR.19) is integral.
Write $\eta=\overline{\xi}$, $b(u)=\overline{g(u)}$,
$w_0(a)=\overline{\varphi(a)}$, and $C=2^e\overline c$.
Then $b(u)^Q=b(u)$ and (BR.10), with $t=2^e$, says
\[
 b(u)-b(u-s)
        =C-\mathbf1_{\langle u\rangle_{2^n}<s}
                              w_0(\langle u\rangle_{2^n}).
 \tag{BR.21}
\]
Also $w_0(1)=1$ and
$w_0(s-1)\eta(-1)=-1$, because $s-1=-1$ modulo $f$.
By (BR.20), $\eta(h)=-1$ in the odd characteristic field $\kappa$.

Suppose the residue of (BR.19) were zero.  Put
\[
 T=\{u\bmod2^n:\eta(u)^{2^e}=1\},\qquad
 \pi:\mathbb Z/2^n\longrightarrow\mathbb Z/2^m.
 \tag{BR.22}
\]
For every unit $y$, changing variables in the zero sum gives
$\sum_u\eta(u)b(uy)=0$: multiplication by $\eta(y)\ne0$
recovers the original sum.  BR1 therefore implies
\[
                         \sum_{u\in T}\eta(u)b(uy)=0.
                         \tag{BR.23}
\]
BR1 also gives $\pi(u)=1$ or $-1$ for $u\in T$.
Define $\epsilon(u)=1$ in the first case and $-1$ in the second.
Since $2^e\mid s$, these two cases give respectively
$u(1-s)=u-s$ and $u(1-s)=u+s$ modulo $2^n$.
The element $1-s$ is a unit.

Use (BR.21) to subtract the two instances $y=1$ and $y=1-s$
of (BR.23).  For each $u\in T$ the resulting summand is
\[
 \eta(u)b(u)-\eta(u)b(u(1-s))
     =\epsilon(u)\eta(u)C-\mathbf1_{u=1}-\mathbf1_{u=-1}.
 \tag{BR.24}
\]
Here are the endpoint checks.  In the case $\pi(u)=1$,
a representative less than $s\le2^m$ must equal $1$,
giving $-\eta(1)w_0(1)=-1$.
In the case $\pi(u)=-1$, apply (BR.21) at $u+s$ and change
the sign.  A representative of $u+s$ less than $s$ must
equal $s-1$, hence $u=-1$; its contribution is
$\eta(-1)w_0(s-1)=-1$.
The two points are distinct since $m\ge2$.

Multiplication by $h=1+2^{n-1}$ permutes $T$, fixes $\pi(u)$,
and multiplies $\eta(u)$ by $-1$.  The first two facts use
$n-1\ge m$ and the evenness of $2^e$.  Hence
\[
                         \sum_{u\in T}\epsilon(u)\eta(u)=0,
                         \tag{BR.25}
\]
since $2$ is invertible in $\kappa$.
Both $1$ and $-1$ belong to $T$.  Summing (BR.24), its left
side is zero by (BR.23), whereas its right side is $-2$ by
(BR.25), a contradiction.
The integral number (BR.19) consequently has nonzero residue,
which is exactly norm one.  This proves (BR.16) with all of
the source's prescribed data.
\end{proof}

\subsection{Unramified homomorphisms and the finite Jacobi relation}

For $N>2$, an odd positive integer $a$ prime to $N$, and
$b\in(\mathbb Z/N)^\times$, set
\[
 \begin{gathered}
 y_b=\langle a^{-1}b\rangle_N,\qquad
 j_{N,a}(b)=\frac{a y_b-\langle b\rangle_N}{N},\\
 E_{N,a}(b)=\frac{a-1}{2}-j_{N,a}(b).
 \end{gathered}
 \tag{BR.26}
\]
These are integers and are exactly the source's \texttt{yCl},
\texttt{jv}, and \texttt{Eone}, evaluated on units.

\begin{proposition}[BR4: the original tame-index condition and Jacobi annihilator]
\label{prop:br-unramified}
Let $\ell$ be odd, let $k$ be any field of characteristic $\ell$,
and let $z$ satisfy (BR.1)--(BR.2).
Then $z$ vanishes on $I_A\cap\ker\gamma$ for every $A$.

Let $F=\mathbb Q(\mu_N)\subset\Omega$ for $N>2$, put
$U=\operatorname{Gal}(\Omega/F)$, and suppose $\gamma|_U=1$.
Choose lifts $\widetilde\sigma_b\in G$ of
$\sigma_b(\zeta_N)=\zeta_N^b$.
For every odd positive $a$ prime to $N$ and every $u\in U$,
\[
 \sum_{b\in(\mathbb Z/N)^\times}
 E_{N,a}(b)\,
 z(\widetilde\sigma_b^{-1}u\widetilde\sigma_b)=0.
 \tag{BR.27}
\]
The statement uses the original $\Omega$ and $k$ and does not require
$k$ finite, a surjective residue map, or a common choice of $H_A$.
\end{proposition}

\begin{proof}
On $K_A=I_A\cap\ker\gamma$, equation (BR.1) makes $z$ an
additive homomorphism.  Its kernel contains $K_A\cap H_A$.
Thus the finite index of that kernel divides
$[K_A:K_A\cap H_A]$, which is prime to $\ell$.
Its image is a finite additive subgroup of $k$ and therefore an
elementary abelian $\ell$-group.  Its order is both a power of
$\ell$ and prime to $\ell$, so it is one.  This proves the first
assertion, without assuming that $H_A$ is contained in inertia.

The restriction $\lambda=z|_U$ is now an additive, locally constant
homomorphism, zero on every $I_A\cap U$.  The group $U$ is compact,
so the locally constant function $\lambda$ has finite image.
Its kernel is open and normal in $U$, and its fixed field is an
actual finite abelian extension of $F$ inside the original $\Omega$.
For every place of $F$, inertia in $G_F$ is exactly
$I_A\cap U$ under this inclusion: both groups fix the chosen
valuation and act trivially on its residue field.
Thus this finite extension is unramified at all finite places.
The field $F$ is totally complex because $N>2$; there is no real
place condition to add.  The finite Hilbert class field theorem
therefore gives an additive map
\[
       \Lambda:\operatorname{Cl}(F)\longrightarrow k,\qquad
       \lambda=\Lambda\circ\operatorname{Art}_F.
       \tag{BR.28}
\]
Only its finite image is used.  The range remains a subgroup of $k$.

We derive the required ideal relation from OS1.
For $1\le r<N$, let
\[
 \theta_{N,r}(t)=
    \left\lfloor\frac{r\langle t\rangle_N}{N}\right\rfloor
    -\left\lfloor\frac{r(N-\langle t\rangle_N)}{N}\right\rfloor.
 \tag{BR.29}
\]
For every nonzero ideal $\mathfrak a$, OS1 constructs an element
whose principal ideal is
\[
                       \prod_t
                 \sigma_t^{-1}\mathfrak a^{\,\theta_{N,r}(t)}.
\]
Consequently the group-ring element
$\sum_t\theta_{N,r}(t)\sigma_t^{-1}$
annihilates the ideal class group.
This consequence uses only the ideal assertion in OS1.

There is no restriction $a<N$ in (BR.27).
Write $a=qN+r$ with $1\le r<N$; this is possible since $(a,N)=1$.
The same floor expression at the positive integer $a$ satisfies
\[
 \theta_{N,a}(t)=\theta_{N,r}(t)+q(2\langle t\rangle_N-N),
 \qquad
 \theta_{N,N-1}(t)=2\langle t\rangle_N-N.
 \tag{BR.30}
\]
Both follow by separating the integer part $q\langle t\rangle_N$.
The relation from OS1 for $r$ and for $N-1$ therefore proves
that $\Theta_a=\sum_t\theta_{N,a}(t)\sigma_t^{-1}$
annihilates $\operatorname{Cl}(F)$ for this arbitrary $a$.

Because $a\langle t\rangle_N$ is not divisible by $N$,
\[
 \theta_{N,a}(t)=
       2\left\lfloor\frac{a\langle t\rangle_N}{N}\right\rfloor
                  -(a-1).
 \tag{BR.31}
\]
Putting $b=at$ modulo $N$ in (BR.26) gives the exact group-ring
identity
\[
 2\sum_b E_{N,a}(b)\sigma_b^{-1}
                       =-\sigma_a^{-1}\Theta_a.
 \tag{BR.32}
\]
Thus the left side annihilates every class.
After applying the additive map $\Lambda$ to it, division by
$2$ is legitimate in the original odd-characteristic field $k$.

Finally Artin is compatible with conjugation:
\[
 \operatorname{Art}_F(\widetilde\sigma_b^{-1}u
                         \widetilde\sigma_b)
                   =\sigma_b^{-1}\operatorname{Art}_F(u).
 \tag{BR.33}
\]
One can check the convention on a Frobenius at a finite prime:
conjugating the automorphism conjugates that prime by
$\sigma_b^{-1}$, and Artin sends its Frobenius to its ideal class.
The resulting compatibility on all elements follows from the
Artin map, or from the density of Frobenius in each finite abelian
quotient.  This is finite reciprocity on the same field $F$;
the inverse agrees with the ideal action in OS1.
Apply $\Lambda$ to (BR.32) at
$\operatorname{Art}_F(u)$ and use (BR.28)--(BR.33).
This proves (BR.27).
\end{proof}

\subsection{The exact Bernoulli scalar on the original coefficient ring}

\begin{proposition}[BR5: the nonzero residual annihilator]
\label{prop:br-scalar}
Let $O=c.O\subset\overline{\mathbb Q}_\ell$ be the source's finite
discrete valuation ring over $\mathbb Z_\ell$, let $k$ be any field
of characteristic $\ell$, and let $\operatorname{red}:O\to k$
be the specified unital ring map.
Let $\chi$ be a nontrivial Dirichlet character of level $N>0$
and order prime to $\ell$, where $\ell$ is odd.
Write $\gamma=\operatorname{redChar}(\chi)$.
Suppose $\gamma$ differs somewhere from the reduction
$\overline{\chi_\ell}$ of the cyclotomic character, and suppose
\[
                         |B_{1,\chi^{-1}}|_\ell=1.
                         \tag{BR.34}
\]
Then $N>2$, $\gamma$ is nontrivial, and there is an odd positive
integer $a$ prime to $N$ such that
\[
 \operatorname{red}\left(
      \sum_{b\in(\mathbb Z/N)^\times}E_{N,a}(b)\chi^{-1}(b)
                    \right)\ne0.
 \tag{BR.35}
\]
The Bernoulli number in (BR.34) is evaluated through the original
inclusion $O\subset\overline{\mathbb Q}_\ell$.
\end{proposition}

\begin{proof}
We first record precisely what the given residue map implies.
The structure map $\mathbb Z_\ell\to O$ is injective, because $O$
is an actual subalgebra of $\overline{\mathbb Q}_\ell$.
The prime ideal $\ker(\operatorname{red})$ contains the nonzero
element $\ell$.  In a discrete valuation ring the only nonzero
prime ideal is the maximal ideal.  Hence this kernel is exactly
the maximal ideal, even though the map need not be surjective.
In particular $\operatorname{red}(x)\ne0$ if and only if $x\in O^\times$.

For use with the norm, all elements of $O$ have norm at most one:
they are integral over $\mathbb Z_\ell$ by module finiteness,
and a monic integral equation contradicts a norm greater than one
by the ultrametric inequality.  A unit and its inverse are integral,
so every unit has norm one.  Write
$\ell=u\varpi^h$ with $u$ a unit and $h>0$, which is possible
because $\ell$ is a nonzero nonunit.
Then $|\varpi|_\ell=\ell^{-1/h}<1$.
The discrete valuation factorization of every nonzero $x\in O$
shows conversely that $|x|_\ell=1$ implies $x$ is a unit.
This proves the norm-to-residue test on the actual ring $O$.

Reduction is injective on roots of unity of order prime to $\ell$.
Indeed, if $x^r=1$, $x\ne1$, and $\ell\nmid r$, then
$1+x+\cdots+x^{r-1}=0$.  Were $\operatorname{red}(x)=1$,
reduction would give $r=0$ in $k$, a contradiction.
Apply this with $r=\operatorname{ord}(\chi)$ to a unit residue
where $\chi$ is not one.  The cyclotomic restriction map
$G\to(\mathbb Z/N)^\times$ is surjective, so $\gamma\ne1$.
Also $N>2$, since the unit groups at levels $1$ and $2$ are trivial.

We choose $a$ so that
\[
                    \operatorname{red}(\chi(a))\ne \overline a.
                    \tag{BR.36}
\]
If $\ell\nmid N$, take $a=\ell$.
It is odd and prime to $N$; the left side is a nonzero character
value whereas the right side is zero.
If $\ell\mid N$, choose the given $g\in G$ where
$\gamma(g)\ne\overline{\chi_\ell}(g)$ and put
$b=\operatorname{cls}_N(g)$.
Choose a positive odd representative $a$ of $b$ modulo $N$.
This is possible: for even $N$ every unit representative is odd;
for odd $N$, adding $N$ changes parity.
The action of $g$ on $\mu_\ell\subset\mu_N$ is the $a$th-power
map.  Thus $\chi_\ell(g)-a\in\ell\mathbb Z_\ell$, and reduction
in the given $k$ gives
$\overline{\chi_\ell}(g)=\overline a$.
The chosen inequality is now (BR.36).

For any nontrivial character $\psi$ with values in a
characteristic-zero field, direct summation in (BR.26) gives
\[
 \sum_b E_{N,a}(b)\psi(b)
                    =(1-a\psi(a))B_{1,\psi}.
                    \tag{BR.37}
\]
Here is the calculation, including the inverse in $y_b$.
The constant term vanishes because $\sum_b\psi(b)=0$.
Multiplication by $a$ permutes the units, so
\[
 \sum_b y_b\psi(b)
      =\psi(a)\sum_t\langle t\rangle_N\psi(t).
 \tag{BR.38}
\]
Multiply the left side of (BR.37) by $N$ and use
$Nj_{N,a}(b)=a y_b-\langle b\rangle_N$.
Equations (BR.38) and (BR.4) give
$N(1-a\psi(a))B_{1,\psi}$, proving (BR.37) since $N\ne0$
in characteristic zero.

Apply (BR.37) to $\psi=\chi^{-1}$ in the specified
$\overline{\mathbb Q}_\ell$.
By (BR.36) the element
$v=1-a\chi^{-1}(a)\in O$ has nonzero reduction and hence is a unit.
Let $S=\sum_b E_{N,a}(b)\chi^{-1}(b)\in O$.
The element $v^{-1}S\in O$ maps to $B_{1,\chi^{-1}}$.
By (BR.34) and the norm criterion above it is a unit.
Thus $S$ is a unit, and its reduction is nonzero.
This proves (BR.35) without making an integrality or nonvanishing
assumption on a different character or coefficient ring.
\end{proof}

\begin{proposition}[BR6: the full residual unramified cocycle theorem]
\label{prop:br-residual}
Under exactly the hypotheses of BR5, every cocycle $z:G\to k$
for $\gamma=\operatorname{redChar}(\chi)$ satisfying
$\operatorname{UnrEv}(\ell,\gamma,z)$ has the form
\[
                     z(g)=(\gamma(g)-1)m\qquad(g\in G)
                     \tag{BR.39}
\]
for some $m$ in the original field $k$.
This is the literal source interface
\texttt{UnramifiedCocycleIs\allowbreak CoboundaryOfBernoulliUnit}.
\end{proposition}

\begin{proof}
Put $F=\mathbb Q(\mu_N)\subset\Omega$ and $U=G_F$.
The character $\gamma$ is trivial on $U$.
For $u\in U$ and any lift $\widetilde\sigma_b$ as in BR4,
the cocycle identity gives
\[
 z(\widetilde\sigma_b^{-1}u\widetilde\sigma_b)
                    =\gamma(\widetilde\sigma_b)^{-1}z(u)
                    =\operatorname{red}(\chi^{-1}(b))z(u).
 \tag{BR.40}
\]
For example, compare $z(u\widetilde\sigma_b)$ with
$z(\widetilde\sigma_b
        (\widetilde\sigma_b^{-1}u\widetilde\sigma_b))$;
the terms $z(\widetilde\sigma_b)$ cancel because $\gamma(u)=1$.
Use in BR4 the integer $a$ supplied by BR5.
Equations (BR.27) and (BR.40) then say
\[
      \operatorname{red}(S)\,z(u)=0,\qquad u\in U.
      \tag{BR.41}
\]
The scalar $\operatorname{red}(S)$ is nonzero, so $z|_U=0$.

Choose $h\in G$ with $\gamma(h)\ne1$, which exists by BR5.
The quotient $G/U=\operatorname{Gal}(F/\mathbb Q)$ is abelian.
Hence $gh$ and $hg$ differ by an element of $U$, on which
both $\gamma-1$ and $z$ vanish.  Thus $z(gh)=z(hg)$.
Using (BR.1) on both products gives
\[
                (\gamma(h)-1)z(g)
                    =(\gamma(g)-1)z(h).
                \tag{BR.42}
\]
Set $m=z(h)/(\gamma(h)-1)\in k$.
Equation (BR.42) proves (BR.39) for all $g$.
In particular no division by the order of
$\operatorname{Gal}(F/\mathbb Q)$ is required; that order may
be divisible by $\ell$.  All local conditions, lifts and the
final coboundary remain on the original $\Omega$ and $k$.
\end{proof}

\subsection{Both characters and the literal CorX endpoint}

\begin{proposition}[BR7: the complete original CorX supplier]
\label{prop:br-cor-x}
Let $\Omega,\ell,c,O,k,\operatorname{red},d,f,n,\varphi,\xi$
be exactly as in the source's \texttt{CorX}: $O=c.O$ is the
specified finite DVR subalgebra, $k$ is any field of characteristic
$\ell$, $\operatorname{red}:O\to k$ is any unital ring map,
and $\operatorname{XData}(\ell,d,\varphi,\xi)$ holds.
Put $\chi=\operatorname{prodX}(\varphi,\xi)$ and let
\[
        \gamma=\operatorname{redChar}(\chi)
          \quad\text{or}\quad
        \gamma=\operatorname{redChar}(\chi)^{-1}.
        \tag{BR.43}
\]
For every locally constant $\gamma$-cocycle $z$ satisfying the
exact condition (BR.2), there is $m\in k$ with
$z(g)=(\gamma(g)-1)m$ for every $g\in G$.
There are no remaining bespoke assumptions in this supplier.
\end{proposition}

\begin{proof}
Write $e=v_2(\ell^d-1)$.  As in BR3,
$\ell$ is odd and $n\ge2e\ge4$.
The order of $\xi$ divides a power of $2$, since it is a
character of the finite group $(\mathbb Z/2^n)^\times$.
The order of $\varphi$ divides $\ell^d-1$.
Therefore
\[
                \chi^{\,(\ell^d-1)2^n}=1,
                \qquad \ell\nmid\operatorname{ord}(\chi).
                \tag{BR.44}
\]
The same statements hold for $\chi^{-1}$.

We verify explicitly both nontriviality and the cyclotomic exclusion
required in BR5.  Put $b=1+2^{n-1}$ modulo $2^n$.
The primitive character calculation (BR.20) gives $\xi(b)=-1$
also in $O$, since its inclusion into
$\overline{\mathbb Q}_\ell$ is injective.
As $\ell f$ is odd, the Chinese remainder theorem gives a unit
$a$ modulo $\ell f2^n$ satisfying
\[
                   a\equiv1\pmod{\ell f},
                   \qquad a\equiv b\pmod{2^n}.
                   \tag{BR.45}
\]
Choose in the original $G$ a lift $g$ of this cyclotomic
automorphism.  It fixes $\mu_\ell$, so
$\overline{\chi_\ell}(g)=1$ in the original $k$.
On the other hand (BR.3) gives
\[
                 \chi(g)=-1,\qquad \chi^{-1}(g)=-1.
                 \tag{BR.46}
\]
Reduction preserves $-1\ne1$ because $\ell$ is odd.
Thus both residual characters are nontrivial and each differs
from $\overline{\chi_\ell}$ at this very same $g$.
In particular $\chi$ and $\chi^{-1}$ are nontrivial before reduction.

The data (BR.5) are preserved under simultaneous inversion:
the conductors and the character orders are unchanged,
$\varphi^{-1}(-1)=\varphi(-1)$ and
$\xi^{-1}(-1)=\xi(-1)$, and the numerical data are identical.
They are also preserved under the actual inclusion
$O\hookrightarrow\overline{\mathbb Q}_\ell$.
For conductor preservation, a character factors through a lower
level exactly when its unit-group kernel contains the corresponding
reduction kernel; an injective coefficient map does not change that
kernel.  Multiplication, inversion and the level changes in
(BR.3) commute with this inclusion.  Therefore BR3 applied to
both pairs gives
\[
                         |B_{1,\chi}|_\ell
                         =|B_{1,\chi^{-1}}|_\ell=1.
                         \tag{BR.47}
\]

For the first alternative of (BR.43), apply BR6 to $\chi$,
using the second equality in (BR.47).  For the second alternative,
apply BR6 to $\chi^{-1}$, using the first equality.
The identity
$\operatorname{redChar}(\chi^{-1})
 =\operatorname{redChar}(\chi)^{-1}$
follows by evaluating the unit-valued character and then applying
the same ring map $\operatorname{red}$.
Thus the input cocycle and its condition (BR.2) are literally
the ones required by BR6 in either case.
Its output lies in the originally quantified $k$.
This proves the full \texttt{CorX} statement.
\end{proof}

\paragraph{Closure and the downstream boundary.}
The complete chain established here is the finite Frobenius projection
and integral partial sum, followed by the prescribed Bernoulli unit;
the original tame-index condition, finite Hilbert class-field descent
and OS1 Jacobi relation; and finally the nonzero scalar, cocycle
vanishing and both-character \texttt{CorX} endpoint.
BR4 uses only the odd-characteristic specialization of the source's
general \texttt{homKill}; it does not assert the unrestricted
integral \texttt{clKill} or arbitrary-coefficient \texttt{homKill}
interfaces.  BR7 supplies the exact input to the paired
\texttt{CountII}/\texttt{CountIII} argument.  Automorphy, rigidity,
congruence-number bounds and the other hypotheses of that argument
are separate obligations.
